% !TeX spellcheck = ru_RU
% !TEX root = vkr.tex

\section{Постановка задачи}
\label{sec:task}

Целью работы является валидация спецификации RISC-V ISA в LLVM относительно Sail.
Для её выполнения были поставлены следующие задачи:
\begin{enumerate}
    \item разработать прототип статического анализатора, который будет включать следующие анализы:
          \begin{itemize}
              \item поиск входных и выходных операндов инструкций;
              \item определение, может ли инструкция читать или писать память;
          \end{itemize}
    \item при помощи статического анализатора сравнить следующие спецификации RISC-V ISA:
          \begin{itemize}
              \item LLVM v19 против sail-riscv\footnote{\url{https://github.com/riscv/sail-riscv/tree/05b845c91d1c1db7b361fc8d06e815b54ca0b07a} (дата обращения: \DTMdate{2025-05-17}).};
              \item Расширение \texttt{P} (Version 0.9.11-draft-20211209\footnote{\url{https://tools.cloudbear.ru/docs/riscv-p-extension-0.9.11-20211209.pdf} (дата обращения: \DTMdate{2025-05-17}).}) в LLVM против sail-riscv-p\footnote{\url{https://github.com/umcann123/sail-riscv-p} (дата обращения: \DTMdate{2025-05-17}).}.
                    % \begin{itemize}
                    %     \item дополнительно определить для каждой инструкции какие контрольно-статусные регистры (CSRs) она читает и пишет.
                    % \end{itemize}
          \end{itemize}
          % \item протестировать анализатор.
          % \item реализовать это (раздел~\ref{subsec:task1});
          % \item спроектировать то-то (раздел~\ref{subsec:task2}) наилучшим образом;
          % \item протестировать на том-то (раздел~\ref{subsec:task3}) и обогнать тех-то;
          % \item \sout{изучить язык \OCaml{}} писать тут не надо, так как тут должны быть задачи, выполнение которых можно проверить/оценить прочитав текст или выслушав доклад;
          %       (т.е. Ваши достижения должны быть опровержимы)
          %       \begin{itemize}
          %           \item это может вызвать сомнения по поводу обзора~--- \emph{выполнить обзор} писать можно и нужно, но защищаемым результатом будут не ваши знания, а текст обзора (то есть он должен иметь ценность сам по себе);
          %       \end{itemize}
          % \item обязательна задача на валидацию результата, будь то эксперимент, апробация, внедрение~--- то есть доказательство того, что Вы сделали что-то, нужное пользователю.
          %       Не путайте с валидацией~--- доказательством того, что Вы сделали то, что хотели Вы (например, тесты~--- валидация результата, хорошо, но недостаточно).
\end{enumerate}
